% !TEX root = ../main.tex
\section{Detalles Importantes de Implementación: Restricciones Redundantes y Simetrías}
\label{sec:implementacion}

La implementación en MiniZinc (\texttt{problemas/PlantaEnergia.mzn}) traduce el modelo matemático de la Sección~\ref{sec:modelo} siguiendo estrictamente la separación declarativa entre modelo y datos (\texttt{datos/Datos.dzn} e \texttt{Instancia01.dzn}--\texttt{Instancia10.dzn}), incorporando validación defensiva de parámetros mediante aserciones estáticas (\texttt{assert}) que verifican $np, n, m \ge 1$, $G \in [0, 100]$, $H \in (0, 100]$, costos y demandas no negativas, y capacidades estrictamente positivas.

\subsection{Eliminación de Reificación Superfluas}

Uno de los aspectos de implementación de mayor impacto sobre el tamaño del modelo consiste en evitar la creación de variables booleanas auxiliares para restricciones que pueden expresarse mediante cotas de dominio. En particular, al sustituir la disyunción reificada sobre el cumplimiento total o parcial de la demanda por las desigualdades directas de consistencia de nodo $d^{\min}_{s,j} \le vendido_{s,j} \le d_{s,j}$ y $0 \le prod_{i,j} \le Cap_i$, el compilador MiniZinc intersecta directamente los intervalos iniciales de cada variable antes de invocar al motor de búsqueda, reduciendo el número de variables booleanas intermedias y eliminando ciclos de propagación reificada.

\subsection{Restricciones Redundantes:}
\label{subsec:redundantes}

En un problema de optimización lineal entera resuelto mediante programación por restricciones, los propagadores de desigualdades lineales operan individualmente mediante consistencia de cotas sobre intervalos enteros. Aunque las restricciones primarias $C_1$--$C_4$ caracterizan el conjunto de soluciones factibles, su propagación local por separado sufre de dos debilidades estructurales: (i) permite explorar ramas con sobreproducción estéril ($\sum_i prod_{i,j} > \sum_s vendido_{s,j}$) que solo se descartan al final por aumentos excesivos de costo sin beneficio, y (ii) en la ecuación objetivo $Z = \sum P_s vendido_{s,j} - \sum c_i prod_{i,j}$, la regla de intervalo calcula la cota superior optimista $\overline{Z}$ evaluando simultáneamente $\overline{vendido}_{s,j} = d_{s,j}$ y $\underline{prod}_{i,j} = 0$, como si fuese posible vender la demanda máxima produciendo cero megavatios a costo nulo.

Para eliminar este desacoplamiento de intervalos y podar drásticamente el árbol \emph{Branch-and-Bound}, se dedujeron e implementaron cinco familias de restricciones redundantes ($R_1$--$R_5$):

\begin{enumerate}
    \item \textbf{Cota inferior de producción diaria agregada ($R_1$):}
    Por transitividad entre el piso mínimo por cliente $vendido_{s,j} \ge d^{\min}_{s,j}$ ($C_2$) y la cobertura diaria $\sum_{s \in S} vendido_{s,j} \le \sum_{i \in P} prod_{i,j}$ ($C_3$), la generación total del parque en cada día $j \in D$ debe cubrir al menos la suma de las demandas mínimas obligatorias del día:
    \begin{equation}
        \forall j \in D: \quad \sum_{i \in P} prod_{i,j} \ge \sum_{s \in S} d^{\min}_{s,j}
        \label{eq:r1_piso_diario}
    \end{equation}
    \emph{Justificación de poda:} Cuando la búsqueda ramifica primero sobre las variables de producción $prod_{i,j}$ antes de instanciar $vendido_{s,j}$, la restricción $R_1$ impide explorar asignaciones de generación cuya suma diaria sea inferior a $\sum_s d^{\min}_{s,j}$, detectando el fallo por consistencia de cotas sin necesidad de descender hasta las variables de venta.

    \item \textbf{Equilibrio exacto de flujo diario sin sobreproducción estéril ($R_2$):}
    Dado que toda central posee costo unitario no negativo ($c_i \ge 0$, y $c_i > 0$ en toda central real) y ninguna restricción técnica impone cotas inferiores a $prod_{i,j}$ (pues $C_4$ solo restringe superiormente la producción en ventanas de régimen alto), cualquier solución factible en la que $\sum_{i \in P} prod_{i,j} > \sum_{s \in S} vendido_{s,j}$ para algún día $j$ puede transformarse reduciendo la producción de las centrales activas en el día $j$ hasta igualar $\sum_{s \in S} vendido_{s,j}$, preservando intactas $C_1$--$C_4$ y aumentando estrictamente la ganancia neta $Z$. Por consiguiente, toda solución óptima satisface el balance exacto de potencia y la cota superior de demanda agregada diaria:
    \begin{equation}
        \forall j \in D: \quad \sum_{i \in P} prod_{i,j} = \sum_{s \in S} vendido_{s,j} \quad \land \quad \sum_{i \in P} prod_{i,j} \le \sum_{s \in S} d_{s,j}
        \label{eq:r2_igualdad_balance}
    \end{equation}
    \emph{Justificación de poda:} Elimina la totalidad del subespacio combinatorio de sobreproducción ($\sum_i prod_{i,j} > \sum_s vendido_{s,j}$) y acopla bidireccionalmente las cotas inferiores y superiores de generación con las de venta en cada día. Previniendo la exploración de ramas con sobre-producción. 

    \item \textbf{Cota lineal acumulada en ventanas de régimen alto ($R_3$):}
    En la restricción disyuntiva primaria $C_4$ \eqref{eq:c4_regimen_alto}, un propagador booleano reificado solo fuerza $prod_{i, k+d^*} \le U_i$ cuando los otros $R_i - 1$ días de la ventana ya han elevado su cota inferior estrictamente por encima de $U_i$. Sin embargo, dado que en toda ventana de $R_i$ días consecutivos a lo sumo $R_i - 1$ días pueden alcanzar la capacidad máxima $Cap_i$ y la producción de al menos un día debe ser $\le U_i$, la suma aritmética de la producción de la planta $i$ en cualquier ventana de longitud $R_i$ satisface incondicionalmente la desigualdad lineal:
    \begin{equation}
        \forall i \in P \text{ con } 1 \le R_i \le n, \; \forall k \in \{1, \dots, n - R_i + 1\}: \quad \sum_{d=0}^{R_i - 1} prod_{i, k+d} \le (R_i - 1)\cdot Cap_i + U_i
        \label{eq:r3_ventana_lineal}
    \end{equation}
    \emph{Justificación de poda:} Al ser una desigualdad lineal directa sobre las variables enteras $prod_{i,j}$, su propagador de cotas reduce el límite superior de energía disponible en bloques de días consecutivos desde la raíz del árbol, antes de que las variables booleanas de la disyunción $C_4$ se fijen.

    \item \textbf{Dominancia económica por orden de mérito y saturación del umbral de régimen normal ($R_4$):}
    Dado que el rango de régimen normal $[0, U_i]$ está libre de penalizaciones inter-temporales en $C_4$, una solución óptima nunca desperdicia capacidad disponible en régimen normal de una planta más económica ni vende a pérdida. Definiendo el orden estricto de mérito entre plantas priorizando aquellas de menor coste de producción $p_1 \prec p_2 \iff (c_{p_1} < c_{p_2}) \lor (c_{p_1} = c_{p_2} \land p_1 < p_2)$ y entre clientes por aquellos que pagan un mayor precio por MW $s_1 \prec s_2 \iff (P_{s_1} > P_{s_2}) \lor (P_{s_1} = P_{s_2} \land s_1 < s_2)$, mediante intercambio marginal ($\pm 1\text{ MW}$ en el día $j \in D$) se demuestran las invariantes:
    \begin{align}
        p_1 \prec p_2:\quad (prod_{p_2, j} > 0) &\implies (prod_{p_1, j} \ge U_{p_1}) \label{eq:r4_merito_normal} \\
        p_1 \prec p_2:\quad (prod_{p_1, j} > U_{p_1} \land prod_{p_2, j} > U_{p_2}) &\implies (prod_{p_1, j} = Cap_{p_1}) \label{eq:r4_ambas_altas} \\
        c_{p_1} < c_{p_2}:\quad (prod_{p_1, j} > U_{p_1} \land prod_{p_2, j} > 0) &\implies (prod_{p_1, j} = Cap_{p_1}) \label{eq:r4_barata_alta} \\
        p_1 \neq p_2 \land c_{p_2} \le c_{p_1}:\quad (prod_{p_1, j} > U_{p_1}) &\implies (prod_{p_2, j} \ge U_{p_2}) \label{eq:r4_cara_alta} \\
        P_s > c_i:\quad (vendido_{s,j} < d_{s,j}) &\implies (prod_{i,j} \ge U_i) \label{eq:r4_cliente_rentable} \\
        P_s < c_i:\quad (prod_{i,j} > 0) &\implies (vendido_{s,j} = d^{\min}_{s,j}) \label{eq:r4_cliente_perdida} \\
        s_1 \prec s_2:\quad (vendido_{s_2, j} > d^{\min}_{s_2, j}) &\implies (vendido_{s_1, j} = d_{s_1, j}) \label{eq:r4_prioridad_clientes}
    \end{align}

    De esta forma se garantiza primero la saturación total de producción en régimen normal y luego en régimen alto priorizando la producción en plantas de bajo costo, y luego la venta a más allá de la demanda mínima a clientes con mayor pago por energía.

    \item \textbf{Proyección lineal de costo marginal ($R_5$):}
    Sea $ganancia\_dia_j = \sum_{s \in S} P_s vendido_{s,j} - \sum_{i \in P} c_i prod_{i,j}$ la ganancia neta del día $j \in D$, tal que $ganancia = \sum_{j \in D} ganancia\_dia_j$. Multiplicando la desigualdad primaria de cobertura $\sum_{s \in S} vendido_{s,j} - \sum_{i \in P} prod_{i,j} \le 0$ por el costo unitario $c_q \ge 0$ de cada planta de referencia $q \in P$ y sumándola a $ganancia\_dia_j$, se obtiene la familia de cotas proyectadas de margen neto diario y global ($\forall j \in D, \forall q \in P$):
    \begin{align}
        ganancia\_dia_j &\le \sum_{s \in S} (P_s - c_q)\cdot vendido_{s,j} - \sum_{i \in P} (c_i - c_q)\cdot prod_{i,j} \label{eq:r5_fm_dia} \\
        ganancia &\le \sum_{j \in D}\sum_{s \in S} (P_s - c_q)\cdot vendido_{s,j} - \sum_{j \in D}\sum_{i \in P} (c_i - c_q)\cdot prod_{i,j} \label{eq:r5_fm_global}
    \end{align}
    Observe que en la ecuación \eqref{eq:r5_fm_dia}, el costo base $c_q$ queda absorbido directamente en el margen neto $(P_s - c_q)$ de cada megavatio vendido. Para $c_q = c_{\min}$, todos los coeficientes $-(c_i - c_{\min}) \le 0$ son no positivos, por lo que el propagador de intervalos deduce instantáneamente en el nodo raíz que $\overline{ganancia\_dia}_j \le \sum_s (P_s - c_{\min}) d_{s,j}$ (descontando de inmediato el costo mínimo inevitable de producir la energía vendida). Más aún, para un escalón intermedio $c_q > c_{\min}$, las plantas más baratas ($c_i < c_q$) adquieren coeficiente positivo $(c_q - c_i) > 0$ y son evaluadas por el propagador en su cota superior $\overline{prod}_{i,j}$ (restringida por $Cap_i$ y $R_3$), mientras que los clientes deficitarios ($P_s < c_q$) adquieren coeficiente negativo $(P_s - c_q) < 0$ y son evaluados en su piso obligatorio $d^{\min}_{s,j}$, proporcionando una cota superior dual exacta que poda subárboles enteros de \emph{Branch-and-Bound}.
\end{enumerate}

\subsection{Rompimiento de Simetrías y Recuperación de Soluciones}
\label{subsec:simetrias}

Al interior de un problema de optimización COP una simetría representa un aspecto de la formulación el modelo que resulta en la posibilidad de que múltiples soluciones substancialmente indistintas dejen invariantes tanto el conjunto de restricciones $\mathcal{C}$ como el valor de la función objetivo $f_{\text{obj}}$, de forma que resultan en gasto computacional adicional sin aportar nuevas soluciones relevantes.\\

Debido a eso se implementan restricciones adicional para limitar las simetrías del problema:

\begin{itemize}
    \item \textbf{Inexistencia de simetrías entre entidades heterogéneas:} Cuando dos centrales $p_1, p_2 \in P$ difieren en algún parámetro de su tripla característica $(c_p, Cap_p, R_p)$, permutar sus vectores de producción altera el costo total de producción en la función objetivo ($c_{p_1} \neq c_{p_2}$), tiene la posibilidad de violar la capacidad máxima ($Cap_{p_1} \neq Cap_{p_2}$) o modifica la factibilidad de las ventanas de régimen alto ($R_{p_1} \neq R_{p_2}$). Análogamente, dos clientes con distinta tarifa $P_{s_1} \neq P_{s_2}$ o distinto vector de demanda no son intercambiables, puesto que alteran el valor de la función objetivo. 

    \item \textbf{Simetrías entre centrales gemelas ($S_1$) y clientes gemelos ($S_2$):} En un modelo genérico donde pueden incorporarse múltiples unidades de la misma tecnología, si dos centrales $p_1 < p_2 \in P$ comparten idéntico costo, capacidad y límite de días en régimen alto, $(c_{p_1}, Cap_{p_1}, R_{p_1}) = (c_{p_2}, Cap_{p_2}, R_{p_2})$, intercambiar sus filas completas de producción $(prod_{p_1, 1..n} \leftrightarrow prod_{p_2, 1..n})$ produce una solución distinta con idéntica ganancia neta y factibilidad. Para romper esta simetría de índice entre plantas gemelas y entre clientes gemelos se imponen los ordenamientos lexicográficos globales:
    \begin{align}
        (c_{p_1}, Cap_{p_1}, R_{p_1}) = (c_{p_2}, Cap_{p_2}, R_{p_2}) \implies \text{lex\_greatereq}\big(prod_{p_1, 1..n},\; prod_{p_2, 1..n}\big) \label{eq:s1_lex_plantas} \\
        P_{s_1} = P_{s_2} \land d_{s_1, 1..n} = d_{s_2, 1..n} \implies \text{lex\_greatereq}\big(vendido_{s_1, 1..n},\; vendido_{s_2, 1..n}\big) \label{eq:s2_lex_clientes}
    \end{align}
    Imponiendo así un orden obligatorio en la solución y evitando el gasto computacional en la búsqueda de soluciones permutadas.\\
    
    Adicionalmente, dentro de cada día $j \in D$, las relaciones de orden de mérito con desempate por índice $(c_{p_1} = c_{p_2} \land p_1 < p_2)$ en \eqref{eq:r4_merito_normal}--\eqref{eq:r4_ambas_altas} rompen la \emph{simetría aditiva intra-día} (que de otro modo permitiría repartir una carga conjunta $E \le U_{p_1} + U_{p_2}$ en cientos de combinaciones equivalentes $(a+k, b-k)$ sin alterar el costo ni el régimen alto), obligando a la asignación canónica de la carga normal a la primera unidad gemela hasta su umbral $U_{p_1}$ antes de cargar la segunda.
    
\end{itemize}
